\section{Random Forest}


\subsection{Random Forest -- From Scratch Implementation}

A \textbf{Random Forest classifier was implemented from scratch} in order to
better understand the internal mechanisms of ensemble learning and to analyze
the impact of different resampling strategies on an imbalanced dataset.

The evaluation was conducted on \textbf{10 resampled datasets}, using fixed
hyperparameters to ensure a fair comparison between resampling methods.

\subsubsection{Experimental Setup}

All experiments were conducted using fixed hyperparameters: \texttt{n\_estimators}=100, \texttt{max\_depth}=10, \texttt{min\_samples\_split}=2, \texttt{min\_samples\_leaf}=1, bootstrap sampling enabled, and all features considered at each split.

% The test set contains \textbf{3477 samples}, with a highly imbalanced class
% distribution:
% \[
% \text{Class 0: } 3275 \quad\quad \text{Class 1: } 202
% \]

% To address this imbalance, three resampling strategies were evaluated:
% \begin{itemize}
%     \item SMOTE
%     \item Tomek Links
%     \item SMOTE + Tomek Links
% \end{itemize}

\subsubsection{Performance Comparison of Resampling Methods}

The table below summarizes the performance of the
from-scratch Random Forest on the test set.

\begin{table}[H]
\centering
\caption{Random Forest (From Scratch) -- Test Performance Comparison}
\label{tab:rf_scratch_comparison}
\begin{tabular}{lccccc}
\toprule
\textbf{Method} & \textbf{Accuracy} & \textbf{Precision} & \textbf{Recall} & \textbf{F1-Score} & \textbf{Time (s)} \\
\midrule
SMOTE & 0.8896 & 0.9447 & 0.8896 & 0.9099 & 5315.15 \\
Tomek Links & \textbf{0.9517} & 0.9439 & \textbf{0.9517} & \textbf{0.9455} & \textbf{2283.09} \\
SMOTE + Tomek Links & 0.8916 & 0.9432 & 0.8916 & 0.9109 & 4948.44 \\
\bottomrule
\end{tabular}
\end{table}





% \subsubsection{Conclusion}

% Based on both quantitative metrics and graphical analysis, \textbf{Tomek Links}
% emerges as the most effective resampling strategy for the from-scratch Random
% Forest implementation. It achieves the highest F1-score and accuracy while
% maintaining the lowest training time and minimal overfitting.

% These results provide a strong baseline for comparison with the
% \texttt{scikit-learn} Random Forest implementation, which is presented in the
% following section.

%%%%%%%%%%%%%%%%%
\subsection{Scikit-learn Implementation}
Training was performed using default hyperparameters unless specified otherwise, and performance was assessed using validation accuracy and class-specific metrics.

\subsubsection{Performance Across Resampling Strategies}

\begin{table}[H]
\centering
\caption{Scikit-Learn Random Forest Global Performance Across Resampling Strategies}
% \label{tab:rf_global_performance}
\begin{tabular}{|l|c|c|c|c|}
\hline
\textbf{Resampling} & \textbf{Training Time (s)} & \textbf{Training Acc.} & \textbf{Validation Acc.} \\ \hline
SMOTE & 4.61 & 0.9599 & 0.9008 \\ \hline
Tomek & 1.29 & 0.9745 & 0.9471 \\ \hline
SMOTE-Tomek & 4.11 & 0.9614 & 0.8973 \\ \hline
\end{tabular}
\end{table}


Similar to the from-scratch implementation, Tomek Links outperformed other resampling methods, achieving the highest validation accuracy (0.9471) and strong class-wise metrics.
